\subsection{Implementierung}

Die Implementierung in Python gliedert sich in sechs Module: \texttt{models.py} (Konstanten, Datenklassen, Fehlertyp \texttt{PlanError}), \texttt{config.py} (Laden und Validieren), \texttt{planner.py} (Bedarf, VLAN-Vergabe, Zuteilung), \texttt{render.py} (Tabelle, CSV, JSON), \texttt{firewall.py} (Regeln) und \texttt{cli.py} (Kommandozeile). Sie verwendet nur die Standardbibliothek (\texttt{ipaddress}, \texttt{math}, \texttt{csv}, \texttt{json}) und \texttt{PyYAML}. Alle Datenklassen sind unveränderlich (\texttt{frozen=True}), und alle Fehler werden als \texttt{PlanError} gemeldet, die die Kommandozeile mit dem Rückgabewert $1$ und der Meldung \glqq Fehler: \dots\grqq{} auf der Fehlerausgabe beantwortet.

\subsubsection{Bedarf und Präfixlänge}

\begin{lstlisting}[language=Python, caption=Bedarf und Präfixlänge (\texttt{planner.py})]
def required_addresses(group: Group, growth_factor: float, gateway_position: str, version: int) -> int:
    """Adressen inkl. Wachstumsreserve und Overhead (Netz, Broadcast, Gateway)."""
    wanted = math.ceil(round(group.hosts * growth_factor, 9))
    overhead = 2 if version == 4 else 1  # IPv6 kennt keinen Broadcast
    if gateway_position != "none":
        overhead += 1
    return wanted + overhead


def prefix_for(total: int, max_prefix: int, name: str) -> int:
    """Kleinste Präfixlänge, deren Block mindestens ``total`` Adressen fasst."""
    prefix = max_prefix - (total - 1).bit_length()
    if prefix < 0:
        raise PlanError(f"Gruppe {name!r}: {total} Adressen passen in keinen Adressraum")
    return prefix
\end{lstlisting}

Die Funktion \code{required_addresses} berechnet $r_i = \lceil\gamma h_i\rceil + o$ nach Definition~\ref{def:bedarf}; das Runden auf neun Nachkommastellen verhindert den Gleitkommafehler aus Lemma~\ref{lem:float}. Die Funktion \code{prefix_for} berechnet $p = w - \lceil\log_2 r\rceil$ mit der Bitlänge aus Lemma~\ref{lem:bitlen} und erfüllt Satz~\ref{satz:minimal}.

\begin{beispiel}
	Für die Gruppe \texttt{Entwicklung} ist $h=120$, $\gamma = 1{,}25$, $o = 3$. Es ist $r = 150 + 3 = 153$, $(r-1).\mathrm{bit\_length}() = 152.\mathrm{bit\_length}() = 8$, denn $152 = 10011000_2$ hat acht Stellen. Also $p = 32 - 8 = 24$.
\end{beispiel}

\subsubsection{VLAN-Vergabe}

\begin{lstlisting}[language=Python, caption=Automatische VLAN-IDs (\texttt{planner.py})]
def assign_vlans(groups: Sequence[Group]) -> Dict[str, int]:
    """Explizite VLAN-IDs übernehmen, fehlende automatisch vergeben (10, 20, ...)."""
    used = {g.vlan_id for g in groups if g.vlan_id is not None}
    result: Dict[str, int] = {}
    candidate = AUTO_VLAN_START
    for g in groups:
        if g.vlan_id is not None:
            result[g.name] = g.vlan_id
            continue
        while candidate in used:
            candidate += AUTO_VLAN_STEP
        if candidate > VLAN_MAX:
            raise PlanError("Keine freie VLAN-ID mehr verfügbar")
        used.add(candidate)
        result[g.name] = candidate
    return result
\end{lstlisting}

Die Menge \code{used} enthält von Anfang an alle expliziten IDs. Dadurch kann eine automatisch vergebene ID nie mit einer expliziten kollidieren, auch wenn diese erst später in der Liste steht (Lemma~\ref{lem:vlan}).

\subsubsection{Zuteilung der Subnetze}

\begin{lstlisting}[language=Python, caption=Kern der Zuteilung (\texttt{planner.py}, Auszug)]
    if cfg.allocation_strategy == "fixed":
        uniform = min(prefix for _, prefix in sized)  # kleinste Länge = größter Block
        sized = [(g, uniform) for g, _ in sized]
        order = sorted(sized, key=lambda item: item[0].priority)
    else:
        order = sorted(sized, key=lambda item: (item[1], item[0].priority))

    vlans = assign_vlans(cfg.groups)
    net_cls = type(base)
    cursor = int(base.network_address)
    end = int(base.broadcast_address)
    subnets: List[Subnet] = []

    for group, prefix in order:
        if prefix < base.prefixlen:
            raise PlanError(
                f"Platzmangel: Gruppe {group.name!r} benötigt ein /{prefix}, "
                f"das Basisnetz {base} ist kleiner"
            )
        size = 1 << (max_prefix - prefix)
        start = -(-cursor // size) * size  # auf Blockgrenze aufrunden
        if start + size - 1 > end:
            raise PlanError(
                f"Platzmangel: Gruppe {group.name!r} (/{prefix}) passt nicht mehr in {base}"
            )
        net = net_cls((start, prefix))
        subnets.append(_build_subnet(group, vlans[group.name], net, cfg.gateway_position))
        cursor = start + size
\end{lstlisting}

Der Code entspricht Schritt~4. Python-Sortierung ist stabil, daher entscheidet bei gleichem Schlüssel die Dateireihenfolge. Die Variable \code{cursor} ist der Zeiger $c$. Die Zeile \code{start = -(-cursor // size) * size} ist das Aufrunden aus Lemma~\ref{lem:align}; bei \code{vlsm} ändert sie den Zeiger nie (Satz~\ref{satz:luecke}), sie ist also eine Sicherung und kein Platzverlust. Die beiden \code{raise}-Anweisungen sind genau die Fälle A und B aus dem Beweis von Satz~\ref{satz:existenz}.

\subsubsection{Freie Bereiche}

\begin{lstlisting}[language=Python, caption=Freie Restbereiche (\texttt{planner.py})]
def _free_ranges(base: IPNetwork, nets: Sequence[IPNetwork]) -> Tuple[IPNetwork, ...]:
    addr = type(base.network_address)
    pos = int(base.network_address)
    end = int(base.broadcast_address)
    free: List[IPNetwork] = []
    for net in sorted(nets, key=lambda n: int(n.network_address)):
        start = int(net.network_address)
        if start > pos:
            free.extend(ipaddress.summarize_address_range(addr(pos), addr(start - 1)))
        pos = int(net.broadcast_address) + 1
    if pos <= end:
        free.extend(ipaddress.summarize_address_range(addr(pos), addr(end)))
    return tuple(free)
\end{lstlisting}

Die Funktion läuft in Adressreihenfolge über die belegten Netze und fasst jede Lücke mit \code{summarize_address_range} zu minimal vielen Blöcken zusammen. Für die Pläne des Algorithmus gibt es nur die Lücke hinter dem letzten Netz (Abschnitt~\ref{sec:frei}); die Funktion ist allgemeiner und wird in den Tests auch mit Lücken vor und zwischen Netzen geprüft.

\subsubsection{Firewall-Regeln}

\begin{lstlisting}[language=Python, caption=Zugriffspaare und nftables-Ausgabe (\texttt{firewall.py})]
def _access_pairs(plan: Plan) -> Iterator[Tuple[Subnet, Subnet]]:
    """(Quelle, Ziel)-Paare; Zugriff auf das eigene Netz braucht keine Regel."""
    by_name = {s.group: s for s in plan.subnets}
    for source in plan.subnets:
        for target in source.access:
            if target != source.group:
                yield source, by_name[target]


def render_nftables(plan: Plan) -> str:
    family = "ip" if plan.base_network.version == 4 else "ip6"
    lines = [
        "# Generiert von subnetplan: Standard ist Verwerfen, erlaubt sind nur die Gruppenzugriffe.",
        "table inet subnet_acl",
        "delete table inet subnet_acl",
        "table inet subnet_acl {",
        "    chain forward {",
        "        type filter hook forward priority 0; policy drop;",
        "        ct state established,related accept",
    ]
    for src, dst in _access_pairs(plan):
        lines.append(
            f'        {family} saddr {src.network} {family} daddr {dst.network} '
            f'accept comment "{src.group} -> {dst.group}"'
        )
    lines += ["    }", "}"]
    return "\n".join(lines)
\end{lstlisting}

Der Generator \code{_access_pairs} erzeugt genau die Menge $P$ aus Definition~\ref{def:zugriff} (mit Wiederholungen, falls ein Ziel mehrfach genannt wird). Die ersten beiden Zeilen der \code{nftables}-Ausgabe, \code{table inet subnet_acl} und \code{delete table inet subnet_acl}, bilden einen bewährten Kniff: Die erste Zeile legt die Tabelle an, falls sie fehlt, die zweite löscht sie, sodass die Datei beim wiederholten Einspielen kein Fehler erzeugt und stets den kompletten Stand ersetzt. Die Kette hat die Standardaktion \code{policy drop}, die Regel \code{ct state established,related accept} entspricht der Regel $0$ aus Abschnitt~\ref{sec:firewall}.

\subsubsection{Validierung der Konfiguration}

\begin{lstlisting}[language=Python, caption=Konsistenzprüfung der Gruppen (\texttt{config.py})]
def _check_groups(groups: Sequence[Group]) -> None:
    names = [g.name for g in groups]
    duplicates = sorted({n for n in names if names.count(n) > 1})
    if duplicates:
        raise PlanError(f"Doppelte Gruppennamen: {duplicates}")

    vlans = [g.vlan_id for g in groups if g.vlan_id is not None]
    dup_vlans = sorted({v for v in vlans if vlans.count(v) > 1})
    if dup_vlans:
        raise PlanError(f"Doppelte VLAN-IDs: {dup_vlans}")

    known = set(names)
    for g in groups:
        missing = sorted(set(g.access) - known)
        if missing:
            raise PlanError(f"Gruppe {g.name!r}: access verweist auf unbekannte Gruppen {missing}")
\end{lstlisting}

Diese Funktion stellt sicher, dass die Voraussetzungen von Lemma~\ref{lem:vlan} (paarweise verschiedene explizite VLAN-IDs) und von Satz~\ref{satz:firewall} (jedes Zugriffsziel ist eine existierende Gruppe, sodass \code{by_name[target]} nie fehlschlägt) erfüllt sind. Die Namensprüfung mit dem regulären Ausdruck \texttt{\^{}[\textbackslash w\textbackslash -\ ]+\$} verhindert, dass Sonderzeichen in Firewall-Kommentare oder ACL-Namen gelangen.

\section{Analyse}

\subsection{Laufzeit}

Im Folgenden sei $n$ die Zahl der Gruppen und $a$ die Gesamtzahl der Einträge in allen \code{access}-Listen. Operationen auf ganzen Zahlen der Breite $w\le128$ zählen als Einheitskosten.

\begin{satz}[Laufzeit der Implementierung]
\label{satz:laufzeit}
	Die Planung (Schritte 1 bis 5) hat die Laufzeit $O(n^2 + a)$, die Ausgabe als Tabelle, CSV oder JSON $O(n + a)$, die Firewall-Ausgabe für \code{iptables} und \code{nftables} $O(n + a)$ und für Cisco $O(n\,(n+a))$. Der Kern (Bedarf, Sortierung, Zuteilung, Freibereiche) benötigt $O(n\log n + w)$.
\end{satz}

\begin{proof}
	\textbf{Schritt 1 (Validierung).} Jede Gruppe wird einmal geprüft ($O(n)$ plus $O(|\text{acc}|)$ je Gruppe, zusammen $O(n + a)$). In \code{_check_groups} wird für jedes Element \code{names.count(n)} beziehungsweise \code{vlans.count(v)} aufgerufen, und jeder Aufruf durchläuft die Liste einmal: $O(n)\cdot O(n) = O(n^2)$. Der Zugriffs-Test \code{set(g.access) - known} kostet $O(|\text{acc}|)$ je Gruppe, zusammen $O(a)$. Gesamt $O(n^2 + a)$.

	\textbf{Schritt 2.} Je Gruppe Einheitskosten: $O(n)$.

	\textbf{Schritt 3.} Sortieren: $O(n\log n)$. \code{assign_vlans}: $O(n)$ plus höchstens $409$ Erhöhungen (Lemma~\ref{lem:vlan}), also $O(n)$.

	\textbf{Schritt 4.} Eine Schleife mit $n$ Durchläufen, je $O(1)$ (die Konstruktion eines Subnetzes ist Einheitskosten): $O(n)$.

	\textbf{Schritt 5.} Es gibt höchstens $O(n)$ Sortier- und Schleifenschritte für die Netze und höchstens $w - p_B$ Blöcke (Satz~\ref{satz:frei}), also $O(n\log n + w)$.

	Die Summe ergibt $O(n^2 + a)$ für die Planung und $O(n\log n + w)$ für den Kern.

	\textbf{Ausgabe.} Jedes Subnetz erzeugt eine Zeile mit konstant vielen Feldern, die Zugriffsliste eine Zeichenkette der Länge $O(|\text{acc}|)$: $O(n + a)$.

	\textbf{Firewall.} \code{_access_pairs} durchläuft alle Subnetze und deren Ziele, $O(n + a)$. Bei \code{iptables} und \code{nftables} wird der Generator einmal durchlaufen. Bei Cisco wird er für jedes der $n$ Subnetze erneut vollständig durchlaufen (\code{for src in plan.subnets:} um \code{for source, dst in _access_pairs(plan)}), das sind $n\cdot O(n+a) = O(n(n+a))$.
\end{proof}

\begin{bemerkung}
	Bei \code{access}-Listen ohne Wiederholungen ist $a \le n(n-1)$. Im ungünstigsten Fall ist $a = \Theta(n^2)$, und die Cisco-Ausgabe ist dann kubisch, obwohl ihre Ausgabegröße nur $\Theta(n + a)$ beträgt. Beide quadratischen Stellen (Validierung und Cisco) sind leicht zu beseitigen (Abschnitt~\ref{sec:ausblick}).
\end{bemerkung}

\subsection{Optimalität des Kerns}

Wir zeigen, dass die Sortierung im Kern nicht vermeidbar ist und die Größe der Ausgabe eine weitere untere Schranke ist. Dazu betrachten wir das \emph{Vergleichsmodell}: Die Prioritäten $\pi_i$ dürfen nur durch Vergleiche (kleiner, gleich, größer) benutzt werden.

\begin{lemma}[Untere Schranke für das Sortieren]
\label{lem:sort}
	Jeder Vergleichsalgorithmus, der $n$ paarweise verschiedene Zahlen sortiert, benötigt im schlimmsten Fall mindestens $\lceil\log_2 n!\rceil \ge \frac{n}{2}\log_2\frac{n}{2}$ Vergleiche.
\end{lemma}

\begin{proof}
	Ein Vergleichsalgorithmus entspricht einem Binärbaum; jedes Blatt gehört zu einer Ausgabepermutation. Da alle $n!$ Permutationen vorkommen müssen, hat der Baum mindestens $n!$ Blätter und damit mindestens die Tiefe $\lceil\log_2 n!\rceil$. Es ist $n! \ge (n/2)^{n/2}$, denn mindestens $n/2$ der Faktoren $1,\ldots,n$ sind $\ge n/2$. Also $\log_2 n! \ge \frac{n}{2}\log_2 \frac{n}{2}$.
\end{proof}

\begin{satz}[Optimalität des Kerns]
\label{satz:optimal}
	Im Vergleichsmodell benötigt jeder korrekte Planungsalgorithmus $\Omega(n\log n)$ Vergleiche im schlimmsten Fall, und seine Ausgabe hat die Größe $\Omega(n + |P|)$, erfordert also mindestens diese Zeit. Der Kern des Algorithmus hat die Laufzeit $O(n\log n + w)$, der Algorithmus ist damit bis auf den konstanten Faktor $w$ \emph{optimal}. Mit Hash-Tabellen statt Listenzählung wird auch die Validierung zu $O(n+a)$, und die Gesamtlaufzeit ist $\Theta(n\log n + a)$.
\end{satz}

\begin{proof}
	\emph{Untere Schranke $\Omega(n\log n)$.} Wähle Gruppen mit gleichem $h_i = 1$, die Strategie \code{fixed}, ein ausreichend großes Basisnetz und paarweise verschiedene Prioritäten $\pi_1,\ldots,\pi_n$. Alle Blöcke haben dieselbe Größe, die Reihenfolge der Ausgabe in Adressreihenfolge ist nach Satz~\ref{satz:korrekt} (K5) genau die Reihenfolge der aufsteigenden Prioritäten. Ein Planungsalgorithmus, der diese Ausgabe liefert, sortiert also die Prioritäten. Nach Lemma~\ref{lem:sort} sind $\Omega(n\log n)$ Vergleiche nötig.

	\emph{Ausgabegröße.} Der Plan enthält $n$ Subnetze und, wenn Firewall-Regeln gewünscht sind, $|P|$ Freigaben. Eine Ausgabe der Größe $\Omega(n + |P|)$ erfordert mindestens ebenso viele Schritte.

	\emph{Obere Schranke.} Satz~\ref{satz:laufzeit} liefert $O(n\log n + w)$ für den Kern. Ersetzt man die Listenzählungen in \code{_check_groups} durch ein Wörterbuch (erwartet $O(1)$ je Zugriff), so ist die Validierung $O(n + a)$; die Cisco-Ausgabe kann mit einer Gruppierung der Paare nach Quelle in $O(n+a)$ erzeugt werden. Die Summe ist $O(n\log n + a)$.
\end{proof}

\begin{bemerkung}
	Die Schranke ist modellabhängig. Sind die Prioritäten ganze Zahlen aus $\{0,\ldots,\Pi\}$, so sortiert Zählsortierung in $O(n + \Pi)$, und die untere Schranke $\Omega(n\log n)$ des Vergleichsmodells gilt nicht. Unbedingt (ohne Modellannahme) gilt nur die lineare Schranke durch die Ausgabegröße.
\end{bemerkung}

\subsection{Speicher}

Die Konfiguration benötigt $O(n + a)$ Speicher, der Plan $O(n + \min(w, \popcnt(F)))$, nämlich $n$ Subnetze und höchstens $w - p_B$ freie Blöcke (Satz~\ref{satz:frei}). Bei einem IPv6-Basisnetz \texttt{/48} sind das bis zu $80$ freie Blöcke, bei IPv4 \texttt{/16} höchstens $16$.

\subsection{Alternative: Buddy-System}

Eine Alternative zur sortierten Zuteilung ist ein binärer Baum über dem Adressraum (ein \emph{Buddy-System}): Jede Anforderung sucht den kleinsten freien Block ausreichender Größe, teilt ihn bei Bedarf und gibt Blöcke wieder zusammen frei. Jede Operation kostet $O(w)$, und das Verfahren unterstützt Anlegen und Löschen im laufenden Betrieb. Für die einmalige Planung aller Gruppen ist die Sortierung mit $O(n\log n)$ einfacher und nach Satz~\ref{satz:kompakt} bereits optimal. Das Buddy-System bietet sich für die in Abschnitt~\ref{sec:ausblick} beschriebene inkrementelle Planung an; der Beweis von Satz~\ref{satz:pack} zeigt, dass das Verfahren \glqq größte Anforderung zuerst\grqq{} dabei ohne Kapazitätsverlust arbeitet.

\subsection{Gesamtkorrektheit}

\begin{satz}[Korrektheit]
\label{satz:gesamt}
	Zu jeder gültigen Konfiguration terminiert der Algorithmus und liefert entweder
	\begin{enumerate}[label=(\alph*)]
		\item einen Plan, der (K1) bis (K6) aus Satz~\ref{satz:korrekt} erfüllt, dessen VLAN-IDs paarweise verschieden sind (Lemma~\ref{lem:vlan}), dessen freie Bereiche die minimale Zerlegung von $B \setminus \bigcup_i N_i$ bilden (Satz~\ref{satz:frei}) und dessen Firewall-Regeln Satz~\ref{satz:firewall} erfüllen, oder
		\item eine Fehlermeldung zum Platzmangel, und dann gilt $S' > |B|$ (Satz~\ref{satz:existenz}). Bei \code{vlsm} gibt es in diesem Fall keinen Plan irgendeiner Art.
	\end{enumerate}
\end{satz}

\begin{proof}
	\emph{Terminierung.} Alle Schleifen laufen über endliche Listen, die einzige \code{while}-Schleife in \code{assign_vlans} erhöht den Kandidaten höchstens $409$-mal (Lemma~\ref{lem:vlan} (iv)). \emph{Ergebnis.} Tritt kein Fehler auf, gelten (K1) bis (K6) nach Satz~\ref{satz:korrekt}, die VLAN-Aussage nach Lemma~\ref{lem:vlan}, die Aussage über freie Bereiche nach Satz~\ref{satz:frei} und die Firewall-Aussage nach Satz~\ref{satz:firewall}. Tritt ein Platzmangel-Fehler auf, folgt die Aussage aus Satz~\ref{satz:existenz}.
\end{proof}

\subsection{Experimentelle Überprüfung}
\label{sec:tests}

Die Testsuite des Projekts umfasst $127$ Tests und erreicht laut Bericht $100\,\%$ Zeilen- und Zweigabdeckung für alle Module (\texttt{cli.py}: $54$ Anweisungen, \texttt{config.py}: $96$, \texttt{firewall.py}: $47$, \texttt{models.py}: $68$, \texttt{planner.py}: $84$, \texttt{render.py}: $35$; zusammen $384$ Anweisungen und $108$ Zweige). Einzelne Tests prüfen unmittelbar die Aussagen dieser Arbeit:

\begin{lstlisting}[language=Python, caption=Test zu (K1) bis (K3) und Satz~\ref{satz:frei} (\texttt{test\_planner.py})]
def test_networks_are_aligned_and_do_not_overlap():
    sizes = [3, 200, 17, 60, 9, 1000, 31]
    groups = [{"name": f"G{i}", "hosts": h} for i, h in enumerate(sizes)]
    plan = make_plan(base_network="10.0.0.0/20", groups=groups)
    nets = [s.network for s in plan.subnets]
    for i, net in enumerate(nets):
        assert net.subnet_of(plan.base_network)
        assert all(not net.overlaps(other) for other in nets[i + 1:])
    covered = sum(n.num_addresses for n in nets) + sum(n.num_addresses for n in plan.free_ranges)
    assert covered == plan.total_addresses
\end{lstlisting}

\begin{lstlisting}[language=Python, caption=Test zum binären Rest (Satz~\ref{satz:frei}) (\texttt{test\_planner.py})]
def test_single_group_leaves_binary_remainder_as_free_ranges():
    plan = make_plan(groups=[{"name": "G", "hosts": 10}])  # 13 + 3 = 16 -> /28
    assert [str(n) for n in plan.free_ranges] == [
        "10.0.0.16/28",
        "10.0.0.32/27",
        "10.0.0.64/26",
        "10.0.0.128/25",
    ]
\end{lstlisting}

Der zweite Test ist ein Spezialfall von Satz~\ref{satz:frei}: Im \texttt{/24} sind $16$ Adressen belegt, $F = 240 = 128 + 64 + 32 + 16$, und die freien Blöcke sind genau die Zweierpotenzen der Binärdarstellung in aufsteigender Reihenfolge.

Vollständige Testabdeckung beweist allerdings nicht die Korrektheit: Sie zeigt, dass jede Zeile und jeder Zweig ausgeführt wurde, nicht dass für \emph{alle} Eingaben die richtigen Werte entstehen. Deshalb wurden die Sätze dieser Arbeit zusätzlich mit einem eigenen Zufallstest gegen die unveränderte Implementierung geprüft (Listing~\ref{lst:prop}, gekürzt). Er erzeugte $20\,000$ Konfigurationen (IPv4 und IPv6, beide Strategien, alle Gateway-Positionen, Wachstumsfaktoren $1{,}0$ bis $2{,}0$, $1$ bis $12$ Gruppen). Davon waren $12\,835$ planbar und $7\,165$ endeten mit Platzmangel; in allen Fällen stimmte das Ergebnis mit den Sätzen überein:

\begin{itemize}
	\item Planbarkeit genau dann, wenn $S' \le |B|$ (Satz~\ref{satz:existenz}).
	\item Alle Netze ausgerichtet, in $B$ enthalten, paarweise disjunkt; lückenlose Belegung ab $a_B$ (Sätze~\ref{satz:korrekt}, \ref{satz:luecke}, \ref{satz:kompakt}).
	\item $u = s - o$ (Lemma~\ref{lem:usable}) und $s < 2r$ bei \code{vlsm} (Satz~\ref{satz:verschnitt}).
	\item Die freien Bereiche haben genau die Größen der gesetzten Bits von $F$, in aufsteigender Reihenfolge (Satz~\ref{satz:frei}).
	\item Die Formel $p = w - \lceil\log_2 r\rceil$ für alle $r \in \{1,\ldots,19\,999\}$ (Satz~\ref{satz:minimal}).
\end{itemize}

Die Packbarkeitsbedingung $(\ast)$ aus Satz~\ref{satz:pack} wurde in einem Universum von $2^5 = 32$ Adressen für $4\,000$ zufällige Belegungen und Aufträge gegen eine vollständige Suche (Backtracking) verglichen; beide Verfahren lieferten stets dasselbe Ergebnis. Schließlich wurde Lemma~\ref{lem:float} an $300\,000$ Zufallspaaren $(h,\gamma)$ mit $h \le 2\cdot10^6$ und zwei Nachkommastellen geprüft (keine Abweichung zur exakten Bruchrechnung).

\begin{lstlisting}[language=Python, label={lst:prop}, caption=Zufallstest der Sätze (gekürzt)]
for it in range(20000):
    cfg = rnd_cfg(strategy, gw, v6)             # Zufallskonfiguration
    sizes = [block_size(g, cfg) for g in cfg.groups]
    n = len(sizes)
    total = n * max(sizes) if strategy == "fixed" else sum(sizes)
    feasible = total <= cfg.base_network.num_addresses
    try:
        plan = plan_subnets(cfg)
    except PlanError:
        assert not feasible                      # Satz Existenz
        continue
    assert feasible
    nets = [s.network for s in plan.subnets]
    pos = int(cfg.base_network.network_address)
    for n_ in nets:                              # Lückenlosigkeit
        assert int(n_.network_address) == pos
        pos += n_.num_addresses
    F = cfg.base_network.num_addresses - (pos - int(cfg.base_network.network_address))
    assert sorted(f.num_addresses for f in plan.free_ranges) == \
           [1 << b for b in range(F.bit_length()) if F >> b & 1]
\end{lstlisting}

